\section{Small auxiliary primes from a uniform zero-free strip}
\label{sec:analytic}

It remains to construct the auxiliary primes in \eqref{eq:auxiliary}
with numerical sizes polynomial in $n$ and $\log p$. We use the following
result from the companion paper. Its proof is not part
of this paper; all dependence on that paper is isolated in this input.

\begin{theorem}[\cite{OpenAIPrimitive26}, Theorem~1.2]\label{analytic:hecke}
For every cyclotomic number field $K$ containing the twelfth roots of
unity, and every finite-order Hecke character $\chi$ of $K$, the Hecke
$L$-function $L_K(s,\chi)$ has no zero in
\[
 \Re s>1-\delta,\qquad \delta=10^{-6}.
\]
The principal character is included, with its pole at $s=1$ permitted.
There is no restriction on conductor or imaginary part.
\end{theorem}

The uniform width is what matters for the application: a zero-free region whose
width deteriorates with the field discriminant or the height would not
give the bound below by the same calculation.

\subsection{A uniform smoothed prime-ideal estimate}

Fix a nonzero nonnegative function $\varphi\in C_c^\infty((1,2))$, and put
\[
 \Phi(s)=\int_0^\infty\varphi(t)t^{s-1}\,dt.
\]
In particular $\Phi(1)>0$. For a number field $M$ of degree $d_M$ and
absolute discriminant $D_M$, define
\[
 \Psi_M(x)=\sum_{\mathfrak p}\sum_{j\ge1}
 \log\Nm\mathfrak p\,
 \varphi\bigl((\Nm\mathfrak p)^j/x\bigr),
\]
where $\mathfrak p$ runs over nonzero prime ideals of its ring of integers.

\begin{lemma}\label{analytic:explicit}
If the Dedekind zeta function $\zeta_M(s)$ has no zero in
$\Re s>1-\delta$, for the fixed $\delta$ in
Theorem~\ref{analytic:hecke}, then, uniformly for $x\ge2$,
\begin{equation}\label{eq:smooth-prime}
 \Psi_M(x)=x\Phi(1)+
 O_\varphi\bigl(x^{1-\delta}(\log D_M+d_M)\bigr).
\end{equation}
The implied constant is independent of $M$.
\end{lemma}

\begin{proof}
We record the uniformity in the standard smoothed explicit formula;
the same formula is used in
\cite[Section~9.2, Equation~(9.9)]{OpenAIPrimitive26}.
Mellin inversion on $\Re s=2$ expresses $\Psi_M(x)$ as the integral
of $-(\zeta_M'/\zeta_M)(s)\Phi(s)x^s$. Moving the contour to
$\Re s=-1/2$ gives
\[
 \Psi_M(x)=x\Phi(1)-\sum_\rho x^\rho\Phi(\rho)
       +O_\varphi(\log D_M+d_M),
\]
where $\rho$ ranges over the nontrivial zeros with multiplicities.
The pole at $1$ gives the first term. The possible zero at $0$ has
order at most $d_M$, and the integral on the new line is bounded using
the functional equation and the absolutely convergent logarithmic
derivative on $\Re s=3/2$. These give the stated remainder.

For completeness, the uniform zero-counting estimate needed here is
\[
 \#\{\rho:j\le|\Im\rho|<j+1\}
 \ll \log D_M+d_M\log(j+2),\qquad j\ge0.
\]
An explicit zero-counting estimate with these uniform dependencies is
given in \cite[Corollary~1.2, Equation~(1.3)]{HasanalizadeShenWong22}.
For $j\ge2$, subtracting its counts at heights $j-1$ and $j+2$
gives the displayed estimate; the count at height $3$ handles $j<2$.
Its absolute implied constant does not depend on $M$.
Integration by parts in the Mellin transform
gives $|\Phi(\sigma+it)|\ll_\varphi(1+|t|)^{-3}$ uniformly for
$0\le\sigma\le1$. Consequently
\[
 \sum_\rho |\Phi(\rho)|\ll_\varphi\log D_M+d_M.
\]
The contour argument can be made first at heights avoiding zeros and
then passed to the limit by this absolute convergence. The assumed
strip places every nontrivial zero in $\Re\rho\le1-\delta$.
Multiplication by $x^{1-\delta}$ proves \eqref{eq:smooth-prime}.
\end{proof}

\subsection{Detecting failure of complete splitting}

\begin{proposition}\label{analytic:auxiliary}
There is an absolute constant $c_0$
such that, for every prime $p>B^{200000}$ and every prime $q\le n$,
an auxiliary prime for $(p,q)$ exists with
\[
 \ell\le c_0B^{20000000}.
\]
\end{proposition}

\begin{proof}
Set
\[
 K=\Q(\mu_{12q}),\qquad M=K(p^{1/q}).
\]
The large-characteristic assumption implies $p\nmid12q$. Thus $p$
is unramified in $K$, and $X^q-p$ is Eisenstein at any prime of $K$
above $p$. It follows that $[M:K]=q$. Since $\mu_q\subset K$, the
extension $M/K$ is cyclic. Abelian zeta factorization~\cite[Chapter~VII]{Neukirch99} writes $\zeta_M$
as the product of the $q$ Hecke $L$-functions corresponding to its
characters over $K$. Theorem~\ref{analytic:hecke} therefore supplies the
strip required in Lemma~\ref{analytic:explicit} for both $K$ and $M$.

We have $[K:\Q]\le12q$ and $[M:\Q]\le12q^2$. The cyclotomic
discriminant formula gives
$\log D_K\le[K:\Q]\log(12q)$. The relative discriminant divides
the power-basis discriminant ideal $(q^qp^{q-1})$, so
\[
 \log D_M\le q\log D_K+
 [K:\Q]\bigl(q\log q+(q-1)\log p\bigr).
\]
Since $q,n,L\le B$, both fields satisfy
$\log D+[\,\cdot\,:\Q]\ll B^3$, with an absolute constant.
By Lemma~\ref{analytic:explicit},
\begin{equation}\label{eq:prime-difference}
 \Psi_K(x)-q^{-1}\Psi_M(x)
  =(1-q^{-1})x\Phi(1)+O_\varphi(x^{1-\delta}B^3).
\end{equation}
Take $x=T B^{20000000}$, where $T\ge2$ is an absolute constant chosen
large enough. Relative to the main term, the error is
$O_\varphi(T^{-\delta}B^{-17})$. Hence the left side of
\eqref{eq:prime-difference} is at least $c_1x$ for a fixed $c_1>0$.

We now bound all contributions that could occur without a prime of
the required kind. On the $K$ side, powers $j\ge2$ and prime ideals
of residue degree at least two contribute at most
\[
 O_\varphi\bigl([K:\Q]\sqrt{x}\log^2(2x)\bigr).
\]
Indeed their underlying rational primes are at most $\sqrt{2x}$;
the sum of residue degrees above each prime is at most $[K:\Q]$;
and the number of relevant exponents is $O(\log(2x))$. The
first-power degree-one prime ideals above $2,3,p,q$ contribute
$O_\varphi([K:\Q]\log(2x))$. Both bounds are $o(x)$ uniformly
at the chosen scale, and enlarging the absolute $T$ makes their sum
less than $c_1x$.

Every remaining first-power degree-one prime of $K$ lies over a
rational prime $\ell\equiv1\pmod{12q}$ outside $\{2,3,p,q\}$.
If $p^{(\ell-1)/q}\equiv1\pmod\ell$, the polynomial $X^q-p$
splits into distinct linear factors in its residue field
$\mathbf F_\ell$. The prime ideal then splits completely in $M/K$,
and its first-power contribution to $\Psi_K$ is cancelled exactly
by the $q$ primes above it in $q^{-1}\Psi_M$. All remaining negative
terms may be discarded when taking an upper bound for the difference.
Thus, if no auxiliary prime lay in the support interval $(x,2x)$,
the preceding error contributions would bound that difference by
less than $c_1x$, a contradiction. We may take $c_0=2T$.
\end{proof}

\begin{proof}[Proof of Theorem~\ref{thm:factorization}]
For each prime $q\le n$ in the large-characteristic branch, search
upward for $\ell_q$, testing primality by trial division and testing
\eqref{eq:auxiliary} by modular exponentiation. The search does not
need the value of $c_0$. Proposition~\ref{analytic:auxiliary} bounds
its length. Construct the table and apply
Theorem~\ref{thm:reduction}. The refined estimates in
Section~\ref{sec:complexity} bound the full sequential bit cost by
$O(B^{400000100})$.
Since $B=O((n+1)L)$, the displayed generous exponent follows.
\end{proof}

The argument also identifies a consequence of the algebraic reduction
alone. GRH for finite-order Hecke $L$-functions places their nontrivial
zeros on $\Re s=1/2$, and therefore implies the strip used above.
Thus the reduction yields deterministic polynomial-time factorization
under GRH independently of the companion's proof of
Theorem~\ref{analytic:hecke}. Theorem~\ref{thm:factorization} uses the
cited theorem to supply the same input without that assumption.
